\section{Definitions and preliminary results}\label{sec:preliminaries}
\subsection{Smooth data and solution classes}
We use the equation, domains, and norms introduced in
Section~\ref{sec:intro}.
All fields are real, and forces need not be divergence free. We write
$L^2_\sigma(D)$ for the divergence-free subspace of $L^2(D;\R^3)$.
The allowed initial velocities and forces are
\begin{align}
 \XX=\mathcal X_{\TT}&=C^\infty_{\mathrm{div}}(\TT;\R^3),&
 \FF=\mathcal F_{\TT}&=C_c^\infty(\TT\times(0,\infty);\R^3),
 \label{eq:inputspaces}\\
 \mathcal X_{\R}&=H^\infty(\R^3;\R^3)\cap L^2_\sigma(\R^3),
 \label{eq:Rinitial}
\end{align}
where $H^\infty=\bigcap_{m=0}^\infty H^m$ has its usual Fr\'echet topology.
On the whole space set
\begin{equation}\label{eq:Rclasses}
 \mathcal F_{\R}=\left\{f\in C^\infty([0,\infty);H^\infty):
 \norm{f}_{L^1_tH^m_x}+\norm{f}_{L^2_tH^m_x}<\infty
 \text{ for every integer }m\ge0\right\}.
\end{equation}
Smoothness into $H^\infty$ means smoothness into each $H^m$, with
one-sided time derivatives at zero. Thus periodic forces vanish near zero
and outside a compact time interval; whole-space forces may be nonzero at
zero and need not have compact support. Both classes consist of forces
defined through any possible singular time of the velocity.

On the torus, velocity and pressure are periodic and $\int_{\TT}p=0$.
On $\R^3$, a classical velocity belongs to $C([0,S];H^m)$ for every
integer $m\ge0$ on each compact interval of its lifespan. Its pressure
gradient is specified below, and the scalar pressure is determined up to
a function of time. In either case local uniqueness defines a maximal
classical lifespan, denoted by $T^\nu_{\max}(a,f)$ on the torus and
$T^\nu_{\max,\R}(a,f)$ on the whole space. A reference solution is
\emph{regular through $T$} if it extends smoothly to $[0,T+\delta]$
for some $\delta>0$ in the relevant class.

For fixed $\nu,T>0$, define
\begin{align}
 \BB_{\nu,a,T}&=\{f\in\FF:T^\nu_{\max}(a,f)\le T\},&
 \BB^0_{\nu,T}&=\BB_{\nu,0,T},\label{eq:singularforces}\\
 \mathcal B^\R_{\nu,a,T}
 &=\{f\in\mathcal F_{\R}:T^\nu_{\max,\R}(a,f)\le T\}.
 \label{eq:Rsingularforces}
\end{align}
We call these the forces producing classical breakdown by $T$.
The inserted solutions satisfy the more specific condition of unbounded
speed at their terminal time.

\subsection{Homogeneous spaces and Fourier multipliers}
Write $\Lambda=(-\Delta)^{1/2}$ and $J=(I-\Delta)^{1/2}$, with the
Fourier conventions of Section~\ref{sec:intro}. On mean-zero periodic
fields, the homogeneous and inhomogeneous Sobolev norms are equivalent
at each fixed order. Appendix~\ref{sec:critical-appendix} specifies the
whole-space homogeneous realization for the positive orders used in the
embeddings.
For the completed energy-force space we fix
\begin{equation}\label{eq:homogeneous-realization}
 \dot H^{-1}(\R^3)=\{h\in\mathcal S':\widehat h=F
 \text{ is a measurable function and }|\xi|^{-1}F\in L^2\}.
\end{equation}
Functions representing distributions are identified almost everywhere.
The map $h\mapsto|\xi|^{-1}\widehat h$ is an isometric bijection onto
$L^2$: its inverse sends $G$ to the inverse transform of $|\xi|G$, which
is tempered because
\[
 \left|\int|\xi|G(\xi)\phi(\xi)\dd\xi\right|
 \le\norm{G}_2\norm{|\xi|\phi}_2\qquad(\phi\in\mathcal S).
\]
Thus this is a separable Hilbert space with no polynomial ambiguity.
The same statement for $H^s$ follows from the isometry
$h\mapsto\langle\xi\rangle^s\widehat h$. Real vector fields correspond
to the closed subspace with $F(-\xi)=\overline{F(\xi)}$.

\subsection{Projection, pressure, and local existence}\label{sec:analytic}
The Leray projection $\PP$ has Fourier symbol
$I-\xi\otimes\xi/|\xi|^2$ on $\R^3\setminus\{0\}$ and the same
symbol at $k\ne0$ on the torus. At the periodic zero mode it is the
identity. It is bounded on every $H^s$ and commutes with derivatives and
the heat semigroup. The projected equation is
\begin{equation}\label{eq:projected}
 \partial_tu-\nu\Delta u=-\PP\nabla\cdot(u\otimes u)+\PP f.
\end{equation}
On the torus, pressure is recovered from
\[
 \Delta p=\nabla\cdot f-\sum_{i,j}\partial_i\partial_j(u_i u_j),
 \qquad \int_{\TT}p=0.
\]
The right side has zero mean. On $\R^3$ we require
\begin{equation}\label{eq:Rpressure}
 \nabla p=(I-\PP)(f-\nabla\cdot(u\otimes u))=:G.
\end{equation}
For smooth $H^\infty$ data, $G$ is smooth and its Fourier transform is
parallel to $\xi$. Hence $\partial_jG_k=\partial_kG_j$, and a potential
is
\[
 p(x,t)=\int_0^1G(rx,t)\cdot x\dd r.
\]
Differentiating and using this symmetry gives
$\partial_jp=\int_0^1\partial_r[rG_j(rx,t)]\dd r=G_j(x,t)$.
There is no requirement that $p\in L^2(\R^3)$. Its gradient belongs to
$L^2$, so a nonzero constant pressure gradient is excluded. On the torus,
the requirement of periodic pressure also excludes an affine pressure.

\begin{proposition}[Local existence, uniqueness, and continuation]
\label{prop:local}\label{lem:Rlocal}
For each initial velocity in the stated class and each force smooth into
every $H^m$ on compact time intervals, equation~\eqref{eq:NS} has a unique
maximal smooth velocity, with pressure determined as above. If a solution
is defined on $[0,S)$, $S<\infty$, and
\begin{equation}\label{eq:criterion}
 \int_0^S\norm{u(t)}_{H^2(D)}^2\dd t<\infty,
\end{equation}
then it extends smoothly beyond $S$.
\end{proposition}
Appendix~\ref{sec:local-appendix} derives this statement from the classical
forced local theory, with one common existence interval for all Sobolev
orders, and proves the displayed continuation criterion. It explains
the viscosity and pressure conventions and the periodic mean reduction;
no compact support of the whole-space velocity is required.

\subsection{Energy and initial vanishing of the compact solution}
We use the fixed solution $(U,P,F)$ chosen after
Theorem~\ref{thm:packet}. Its energy estimate supplies the time-integrated
dissipation bound and the initial vanishing needed for rescaling.

\begin{lemma}[Energy, dissipation and initial vanishing]\label{lem:packetenergy}
The localized solution satisfies
\[
 M:=\sup_{0\le t<1}\norm{U(t)}_2<\infty,\qquad
 D:=\norm{\nabla U}_{L^2((0,1)\times\R^3)}<\infty.
\]
Both $U$ and $P$ vanish on an initial time interval and therefore extend smoothly by zero to negative times.
\end{lemma}
\begin{proof}
On each closed interval $[0,b]$ with $b<1$, compact spatial support justifies the energy identity
\[
 \tfrac12(\norm{U(t)}_2^2)'
 +\nu\norm{\nabla U(t)}_2^2=\ip{F(t)}{U(t)}.
\]
Set $N(t)=\int_0^t\norm{F(s)}_2\dd s$. Dropping dissipation gives
$\tfrac12(\norm{U}_2^2)'\le\norm{F}_2\norm{U}_2$.
For $\zeta>0$, the derivative of $(\norm{U}_2^2+\zeta^2)^{1/2}$ is at most $\norm{F}_2$. Integrating from zero and letting $\zeta\downarrow0$ proves $\norm{U(t)}_2\le N(t)$.

Substituting this bound into the integrated energy identity gives
\begin{equation}\label{eq:packetenergy}
 \norm{U(t)}_2^2+2\nu\int_0^t\norm{\nabla U(s)}_2^2\dd s
 \le 2\int_0^t\norm{F(s)}_2N(s)\dd s=N(t)^2.
\end{equation}
Since $F$ is smooth and time compact, the right side is bounded independently of $t<1$. Monotone convergence of the dissipation integral proves $D<\infty$.

The temporal support of $F$ is a compact subset of $(0,\infty)$, so $F=0$ on $[0,\tau]$ for some $\tau>0$. Estimate~\eqref{eq:packetenergy} gives $U=0$ there. Equation~\eqref{eq:packet} then implies $\nabla P=0$, and compact spatial support fixes this spatial constant to zero. Vanishing on an interval proves smoothness of both zero extensions.
\end{proof}
